\begin{mainquestion}{diodethevenin}{I Love Diodes}

\questionpoints
Consider the circuit below. Using a constant voltage drop model, determine the minimum input voltage ($V_{in}$) required to turn on the LED.

Assume the following component values:
\begin{itemize}
    \item $R_1 = 1\text{k}\Omega$
    \item $R_2 = 1\text{k}\Omega, \quad R_3 = 4\text{k}\Omega$
    \item $R_4 = 4\text{k}\Omega, \quad R_5 = 1\text{k}\Omega$
    \item The LED ($D_1$) has a forward voltage drop $V_f = 1.8\text{V}$.
\end{itemize}

\begin{center}
\begin{circuitikz}[american, scale=1.2] 
    % --- Coordinates (Shortened Vertically) ---
    
    % Input Terminals (Pulled far left)
    \coordinate (InTop) at (-7.5, 2.5);
    \coordinate (InBot) at (-7.5, -2.5);
    
    % Resistor Top/Bot (Kept at original x=-5, but compressed y)
    \coordinate (ResTop) at (-5, 2.5);
    \coordinate (ResBot) at (-5, -2.5);
    
    % Heart Shape (Compressed Vertically)
    \coordinate (TopDip) at (0, 1.0);       % Lowered dip
    \coordinate (LeftLobe) at (-1.5, 2.5);  % Lowered peaks
    \coordinate (RightLobe) at (1.5, 2.5);
    \coordinate (LeftElbow) at (-3, 0.5);   % Elbows closer to center
    \coordinate (RightElbow) at (3, 0.5);
    \coordinate (BottomTip) at (0, -2.5);   % Raised bottom tip

    % --- The Extended "I" (Input Section) ---
    
    % 1. The Vertical Resistor (The "Body" of the I)
    \draw (ResTop) to[R, l=$R_1$] (ResBot);

    % 2. The Horizontal Extensions (The "Serifs" of the I)
    % Top wire extending left to Vin
    \draw (InTop) to[short, o-] (ResTop);
    % Bottom wire extending left to Vin
    \draw (InBot) to[short, o-] (ResBot);
    
    % Labels for Vin at the far left terminals
    \node at (InTop) [above left] {$+$};
    \node at (InBot) [below left] {$-$};
    \node at (-8, 0) {$V_{in}$}; % Label centered between terminals

    % --- Connecting "I" to Heart ---
    % Top wire connects Resistor to Upper Left Lobe
    \draw (ResTop) -- (LeftLobe);
    
    % Bottom wire connects Resistor to Bottom Tip
    \draw (ResBot) -- (BottomTip);

    % --- The Heart Bridge ---
    % The Top Rail: Connect Left Lobe to Right Lobe via the Dip
    % This forces the entire top 'M' shape to be the positive rail
    \draw (LeftLobe) to[short] (TopDip) to[short] (RightLobe);

    % Left Branch
    \draw (LeftLobe) to[R, l_=$R_2$, *-*] (LeftElbow) node[label={left:$A$}]{};
    \draw (LeftElbow) to[R, l_=$R_3$] (BottomTip);

    % Right Branch
    \draw (RightLobe) to[R, l=$R_4$, *-*] (RightElbow) node[label={right:$B$}]{};
    \draw (RightElbow) to[R, l=$R_5$] (BottomTip);

    % The LED Bridge (Center)
    \draw (LeftElbow) to[D, l_=$D_1$] (RightElbow);

\end{circuitikz}
\end{center}

% =========================================================
\begin{partquestion}{theveninvoltage}
{To analyze the LED, it can be helpful to first remove it from the circuit. Calculate the Thevenin Voltage ($V_{th}$) between terminals $A$ and $B$ in terms of $V_{in}$.}

\begin{answer}[2in]
        The voltage at nodes A and B is determined by the voltage divider rule (since the top rail is $V_{in}$ and the bottom rail is ground).
        \[ V_A = V_{in} \left( \frac{R_3}{R_2 + R_3} \right) = V_{in} \left( \frac{4\text{k}}{1\text{k} + 4\text{k}} \right) = 0.8 V_{in} \]
        \[ V_B = V_{in} \left( \frac{R_5}{R_4 + R_5} \right) = V_{in} \left( \frac{1\text{k}}{4\text{k} + 1\text{k}} \right) = 0.2 V_{in} \]
        The Thevenin voltage is the difference between these nodes:
        \[ V_{th} = V_A - V_B = 0.8 V_{in} - 0.2 V_{in} = 0.6 V_{in} \]
\end{answer}

\pagebreak
\end{partquestion}

% =========================================================
\begin{partquestion}{theveninresistance}
{Calculate the Thevenin Resistance ($R_{th}$) looking into terminals $A$ and $B$.}

\begin{answer}[2in]
        To find $R_{th}$, we turn off the source ($V_{in} \rightarrow$ short circuit). This connects the top rail directly to the bottom rail.
        \begin{itemize}
            \item At node A: $R_2$ and $R_3$ are now in parallel.
            \[ R_A = R_2 || R_3 = \frac{1 \cdot 4}{1 + 4} = 0.8\text{k}\Omega \]
            \item At node B: $R_4$ and $R_5$ are now in parallel.
            \[ R_B = R_4 || R_5 = \frac{4 \cdot 1}{4 + 1} = 0.8\text{k}\Omega \]
        \end{itemize}
        The total resistance is the sum of these two branches:
        \[ R_{th} = R_A + R_B = 0.8\text{k}\Omega + 0.8\text{k}\Omega = 1.6\text{k}\Omega \]
\end{answer}

\end{partquestion}

% =========================================================
\begin{partquestion}{mininput}
{Using your Thevenin equivalent circuit, determine the minimum input voltage ($V_{in}$) required to turn on the LED.}

\begin{answer}[2in]
        We can model the heart circuit as a single voltage source $V_{th}$ in series with $R_{th}$ and the LED.
        
        For the LED to turn on, the Thevenin voltage must exceed the forward voltage drop $V_f$:
        \[ V_{th} > V_f \]
        
        Substitute the expressions we found in parts (a) and (b):
        \[ 0.6 V_{in} > 1.8\text{V} \]
        
        Solve for $V_{in}$:
        \[ V_{in} > \frac{1.8\text{V}}{0.6} \]
        \[ V_{in} > 3\text{V} \]
        
        The minimum input voltage required is \textbf{3 Volts}.
\end{answer}

\end{partquestion}

\end{mainquestion}
